\begin{tcolorbox}[colback=red!3!white,colframe=red!50!black,boxrule=0.8pt,title={Korrekturhinweise},fonttitle=\bfseries]
\textbf{Korrekturen aus der Rechenprüfung (30.~Sept.~2026).} Die folgenden Stellen sind im Text korrigiert bzw. vermerkt; jede trägt dort einen datierten Hinweis mit dem vorherigen Wert:
\begin{itemize}
\item $E_0$: Die beiden Werte sind derselbe $E_0$, nackt ($\sqrt{m_e m_\mu} = 7{,}348$~MeV) und fraktal korrigiert ($E_0/\sqrt{K_{\text{frak}}} = 7{,}398$~MeV, Dok.~190, R72). Der korrigierte Wert dient als Anker für $\alpha$; aus den Massen allein folgt $\alpha = 1/137{,}059$ ($+0{,}017\,\%$).
\item $\alpha$: Übereinstimmung $0{,}0005\,\%$ statt $0{,}03\,\%$.
\item Anomale magnetische Momente: Die Kurzformel $a_\ell = \frac{\xi}{2\pi}(E_\ell/E_e)^2$ ist entfernt; sie ergab nachgerechnet $a_\mu = 0{,}907$ statt $1{,}166 \times 10^{-3}$. Maßgeblich sind Dok.~018 und 158.
\item Dimension der Lagrange-Dichte: $[E^4]$, nicht $[E^2]$.
\item $G_{\text{SI}}$-Formel: $\frac{\xi^2}{4m_e}\times C_{\text{dim}}\times C_{\text{conv}} \to \frac{\xi^2}{4m_e}\times C_{\text{conv}}\times K_{\text{frak}}$ -- mit $C_{\text{dim}}$ ohne den Ausgleichsfaktor $E_{\text{char}}$ wäre $G$ um den Faktor $E_{\text{char}} = 28{,}8$ zu klein. Ebenso: $\xi^2/(4m_e)$ hat $[E^{-1}]$, nicht $[E^{-2}]$.
\item Leptonmassen-Tabelle: $0{,}507/103{,}5/1815 \to 0{,}505/104{,}96/1783{,}4$~MeV (Korpusleiter).
\end{itemize}
\end{tcolorbox}

\section*{Abstract}
		Das T0-Modell beschreibt physikalische Phänomene durch ein universelles Energiefeld $E_{\text{field}}(x,t)$ mit dem Parameter $\xi = \frac{4}{3} \times 10^{-4}$. Die Feldgleichung ist $\square E_{\text{field}} = 0$, die Lagrange-Dichte $\mathcal{L} = \xi (\partial E)^2$. 
		
		\textbf{T0-Einheiten:} Alle Konstanten werden auf 1 gesetzt: $c = \hbar = \alpha = G = 1$.
		
		\textbf{Fundamentale Größen:}
		\begin{itemize}
			\item Parameter: $\xi = \frac{4}{3} \times 10^{-4}$ (einziger freier Parameter!)
			\item Charakteristische Energie: $E_0 = \sqrt{m_e \cdot m_\mu} = 7{,}348$ MeV (nackt); fraktal korrigiert $E_0/\sqrt{K_{\text{frak}}} = 7{,}398$ MeV (Anker für $\alpha$)
		\end{itemize}
		
		\textbf{Vorhersagen (SI-Werte):} Leptonmassen mit 2\% Genauigkeit, Feinstrukturkonstante $\alpha \approx 1/137$ (aus den Massen allein $1/137{,}059$, $+0{,}017\,\%$).
		
		\textbf{Detaillierte Herleitungen:} Siehe Dokument 011 (Feinstruktur), 012 (Gravitation), 018 (g-2), 019 (Lagrangian).

	
	
	
	% ============================================================================
	\section{Einheitenkonvention}
	\label{sec:units}
	
	\subsection{T0 natürliche Einheiten}
	
	T0 setzt **alle** fundamentalen Konstanten auf 1:
	
	\begin{equation}
		\boxed{c = \hbar = \alpha = G = 1}
	\end{equation}
	
	\textbf{Konsequenzen:}
	\begin{itemize}
		\item Alle Größen dimensionslos oder in Potenzen einer Einheit
		\item Maximale Vereinfachung der Formeln
		\item Energie = Masse = Länge$^{-1}$ = Zeit$^{-1}$
		\item $4\pi\varepsilon_0 = 1$ (folgt aus $\alpha = e^2 = 1$)
	\end{itemize}
	
	\subsection{Rekonstruktion der SI-Werte}
	
	\textbf{Wichtig:} Obwohl in T0 $\alpha = 1$, kann der SI-Wert rekonstruiert werden!
	
	\textbf{Feinstrukturkonstante (SI):}
	\begin{equation}
		\alpha_{\text{SI}} = \xi \cdot \left(\frac{E_0}{1\,\text{MeV}}\right)^2 \approx \frac{1}{137}
	\end{equation}
	
	mit:
	\begin{itemize}
		\item $\xi = \frac{4}{3} \times 10^{-4}$ (T0-Parameter)
		\item $E_0 = 7{,}398$ MeV (charakteristische Energie)
	\end{itemize}
	
	\textbf{Gravitationskonstante (SI):}
	\begin{equation}
		G_{\text{SI}} = \frac{\xi^2}{4m_e} \times C_{\text{conv}} \times K_{\text{frak}}
	\end{equation}
	\textit{(Korrigiert am 30.~Sept.~2026: vorher $\frac{\xi^2}{4m_e}\times C_{\text{dim}}\times C_{\text{conv}}$ -- $C_{\text{dim}} = 1/E_{\text{char}} = 3{,}472\times10^{-2}$ wird in der Schrittkette von Dok.~012 durch den Faktor $E_{\text{char}} = 28{,}8$ wieder aufgehoben; in der Formel allein ergäbe das Produkt $2{,}35\times10^{-12}$ statt $6{,}674\times10^{-11}$.)}
	
	\textbf{Prinzip:} 
	\begin{itemize}
		\item In T0-Rechnungen: alle Konstanten = 1 (einfache Formeln)
		\item Für Experimente: SI-Werte aus $\xi$ berechnen
		\item Beide äquivalent, nur verschiedene Darstellungen!
	\end{itemize}
	
	\subsection{Dimensionen in T0-Einheiten}
	
	Mit $c = \hbar = \alpha = G = 1$:
	
	\begin{align}
		[E] &= 1 \quad \text{(alle Größen dimensionslos oder Potenzen von Energie)} \\
		[m] &= 1 \\
		[t] &= 1 \\
		[L] &= 1 \\
		[\partial_\mu] &= 1
	\end{align}
	
	Alle physikalischen Größen werden in einer einzigen Einheit gemessen!
	
	% ============================================================================
	\section{Zeit-Energie-Dualität}
	\label{sec:duality}
	
	\subsection{Fundamentale Relation}
	
	\begin{equation}
		T_{\text{field}}(x,t) \cdot E_{\text{field}}(x,t) = 1
		\label{eq:duality}
	\end{equation}
	
	mit $[T_{\text{field}}] = E^{-1}$ und $[E_{\text{field}}] = E$.
	
	\subsection{Intrinsisches Zeitfeld}
	
	\begin{equation}
		T_{\text{field}}(x,t) = \frac{1}{E_{\text{field}}(x,t)}
	\end{equation}
	
	% ============================================================================
	\section{Universelle Feldgleichung}
	\label{sec:field_equation}
	
	\subsection{Wellengleichung}
	
	\begin{equation}
		\square E_{\text{field}} = 0
	\end{equation}
	
	mit d'Alembert-Operator:
	\begin{equation}
		\square = \nabla^2 - \frac{\partial^2}{\partial t^2}
	\end{equation}
	
	\subsection{Mit Quellen}
	
	\begin{equation}
		\nabla^2 E_{\text{field}} = 4\pi G \rho \cdot E_{\text{field}}
	\end{equation}
	
	Dimensionscheck: $[E^3] = [E^{-2}][E^4][E] = [E^3]$ ✓
	
	% ============================================================================
	\section{Lagrange-Dichte}
	\label{sec:lagrangian}
	
	\subsection{Universelle Lagrange-Dichte}
	
	\begin{equation}
		\mathcal{L} = \xi \cdot (\partial_\mu E_{\text{field}})(\partial^\mu E_{\text{field}})
	\end{equation}
	
	mit $\xi = \frac{4}{3} \times 10^{-4}$.
	
	\subsection{Euler-Lagrange-Gleichung}
	
	\begin{equation}
		\frac{\partial \mathcal{L}}{\partial E} - \partial_\mu \frac{\partial \mathcal{L}}{\partial(\partial_\mu E)} = 0
	\end{equation}
	
	ergibt:
	\begin{equation}
		\square E_{\text{field}} = 0
	\end{equation}
	
	% ============================================================================
	\section{Charakteristische Längen}
	\label{sec:characteristic_lengths}
	
	\subsection{T0-charakteristische Länge}
	
	\begin{equation}
		r_0 = 2GE
	\end{equation}
	
	Dimension: $[r_0] = [E^{-2}][E] = [E^{-1}] = [L]$ ✓
	
	\subsection{Herleitung}
	
	Für sphärisch symmetrische Punktquelle $\rho(r) = E_0 \delta^3(\vec{r})$:
	
	Lösung von $\nabla^2 E = 4\pi G \rho E$:
	\begin{equation}
		E(r) = E_0 \left(1 - \frac{r_0}{r}\right)
	\end{equation}
	
	mit $r_0 = 2GE_0$.
	
	\subsection{Zeitskala}
	
	\begin{equation}
		t_0 = \frac{r_0}{c} = r_0 = 2GE
	\end{equation}
	
	(da $c = 1$)
	
	% ============================================================================
	\section{Charakteristische Energie}
	\label{sec:characteristic_energy}
	
	\subsection{Definition}
	
	Die charakteristische Energie $E_0$ ist das geometrische Mittel der Elektron- und Myonmasse (Herleitung in Dokument 011):
	
	\begin{equation}
		\boxed{E_0 = \sqrt{m_e \cdot m_\mu}}
		\label{eq:E0_definition}
	\end{equation}
	
	\subsection{Numerische Werte}
	
	Aus experimentellen Massen:
	\begin{align}
		E_0 &= \sqrt{0{,}511 \times 105{,}66} \\
		&= \sqrt{53{,}99} \\
		&= 7{,}348 \text{ MeV}
	\end{align}
	
	Fraktal korrigierter Wert (Dok.~190, R72):
	\begin{equation}
		E_0^{\text{T0}} = \frac{E_0}{\sqrt{K_{\text{frak}}}} = \frac{7{,}348}{\sqrt{1 - 100\,\xi}} = 7{,}397 \approx 7{,}398 \text{ MeV}
	\end{equation}

	Beide Werte sind derselbe $E_0$, nackt und fraktal korrigiert: Absolutwerte tragen $K_{\text{frak}}$, Verhältnisse nicht. Der Wurzelfaktor folgt daraus, dass $E_0$ quadratisch in $\alpha = \xi E_0^2$ eingeht. Der korrigierte Wert $7{,}398$~MeV dient als Anker für $\alpha$.

	\textit{(Korrigiert am 30.~Sept.~2026: vorher \enquote{Theoretischer T0-Wert: $E_0^{\text{T0}} = 7{,}398$~MeV; Abweichung: 0,7\,\% (im Rahmen geometrischer Korrekturen)}. Der Zusammenhang $7{,}398/7{,}348 = 1{,}00682 \approx 1/\sqrt{K_{\text{frak}}} = 1{,}00673$ ($0{,}008\,\%$) war nicht angegeben.)}

	\subsection{Verwendung}
	
	$E_0$ dient als Energieskala für:
	\begin{itemize}
		\item Feinstrukturkonstante: $\alpha = \xi (E_0/1\,\text{MeV})^2$
		\item Normierung elektromagnetischer Effekte
	\end{itemize}
	
	% ============================================================================
	\section{Der Parameter $\xi$}
	\label{sec:xi_parameter}
	
	\subsection{Definition}
	
	\begin{equation}
		\boxed{\xi = \frac{4}{3} \times 10^{-4} = 1{,}3333 \times 10^{-4}}
	\end{equation}
	
	Dimensionslos: $[\xi] = 1$.
	
	\subsection{Geometrische Komponenten}
	
	\begin{equation}
		\xi = G_3 \times S_{\text{ratio}}
	\end{equation}
	
	wobei:
	\begin{itemize}
		\item $G_3 = \frac{4}{3}$: Geometrischer Faktor (Kugel-Würfel-Verhältnis)
		\item $S_{\text{ratio}} = 10^{-4}$: Skalenverhältnis
	\end{itemize}
	
	% ============================================================================
	\section{Skalenhierarchie}
	\label{sec:scale_hierarchy}
	
	\subsection{Planck-Länge als Referenz}
	
	\begin{equation}
		\ell_P = \sqrt{G} = 1 \quad \text{(in nat. Einheiten)}
	\end{equation}
	
	\subsection{Skalenverhältnis}
	
	\begin{equation}
		\xi_{\text{ratio}} = \frac{\ell_P}{r_0} = \frac{\sqrt{G}}{2GE} = \frac{1}{2\sqrt{G} \cdot E}
	\end{equation}
	
	Für $E \sim 1$ GeV:
	\begin{equation}
		\frac{r_0}{\ell_P} = 2\sqrt{G}\,E = \frac{2E}{E_P} \sim 1{,}6\times10^{-19} \quad \text{(sub-Planck)}
	\end{equation}
	\textit{(Korrigiert am 29.~Sept.~2026: vorher $r_0/\ell_P \sim 10^7$ -- mit $r_0 = 2GE$ und $\ell_P = \sqrt{G}$ gilt $r_0/\ell_P = 2E/E_P$, für $E = 1$~GeV also $1{,}6\times10^{-19}$.)}
	
	% ============================================================================
	\section{Teilchen als Feldanregungen}
	\label{sec:particles}
	
	\subsection{Klassifikation nach Energie}
	
	\begin{table}[h]
		\centering
		\begin{tabular}{lc}
			\hline
			\textbf{Teilchen} & \textbf{Energie [MeV]} \\
			\hline
			Elektron & 0,511 \\
			Myon & 105,658 \\
			Tau & 1776,86 \\
			\hline
		\end{tabular}
	\end{table}
	
	\subsection{Antiteilchen}
	
	Negative Feldanregungen: $E_{\text{field}} < 0$
	
	% ============================================================================
	\section{Feinstrukturkonstante}
	\label{sec:fine_structure}
	
	\subsection{T0-Herleitung}
	
	Die Feinstrukturkonstante folgt aus $\xi$ und $E_0$ (Herleitung in Dokument 011):
	
	\begin{equation}
		\boxed{\alpha = \xi \cdot \left(\frac{E_0}{1\,\text{MeV}}\right)^2}
		\label{eq:alpha_derivation}
	\end{equation}
	
	\subsection{Numerische Berechnung}
	
	Mit $\xi = \frac{4}{3} \times 10^{-4}$ und $E_0 = 7{,}398$ MeV:
	
	\begin{align}
		\alpha &= 1{,}3333 \times 10^{-4} \times (7{,}398)^2 \\
		&= 1{,}3333 \times 10^{-4} \times 54{,}73 \\
		&= 7{,}297 \times 10^{-3} \\
		&= \frac{1}{137{,}04}
	\end{align}
	
	Experimentell: $\alpha_{\text{exp}} = \frac{1}{137{,}036}$
	
	Übereinstimmung: 0,0005\,\% (mit dem Anker $E_0 = 7{,}398$~MeV). Aus den Massen allein, $\alpha = \xi E_0^2/K_{\text{frak}}$ mit $E_0 = 7{,}348$~MeV, folgt $\alpha = 1/137{,}059$ ($+0{,}017\,\%$).

	\textit{(Korrigiert am 30.~Sept.~2026: vorher \enquote{Übereinstimmung: 0,03\%}; $1/137{,}035$ gegen $1/137{,}036$ sind $0{,}0005\,\%$.)}
	
	\subsection{Dimensionscheck}
	
	\begin{equation}
		[\alpha] = [\xi] \times \left[\frac{E}{E}\right]^2 = 1 \times 1 = 1 \quad \checkmark
	\end{equation}
	
	% ============================================================================
	\section{Gravitationskonstante}
	\label{sec:gravitational_constant}
	
	\subsection{T0-Formel}
	
	Die Gravitationskonstante wird aus $\xi$ und $m_e$ hergeleitet (Herleitung in Dokument 012):
	
	\begin{equation}
		G = \frac{\xi^2}{4m_e} \times C_{\text{conv}} \times K_{\text{frak}}
		\label{eq:G_formula}
	\end{equation}
	\textit{(Korrigiert am 30.~Sept.~2026: wie oben, vorher $\times\,C_{\text{dim}}\times C_{\text{conv}}$ ohne $K_{\text{frak}}$.)}
	
	wobei:
	\begin{itemize}
		\item $C_{\text{conv}} = 7{,}783\times10^{-3}$: SI-Umrechnungsfaktor
		\item $K_{\text{frak}} = 0{,}986$: fraktale Korrektur
	\end{itemize}
	
	\subsection{Fundamentale Beziehung}
	
	In natürlichen Einheiten:
	\begin{equation}
		\xi = 2\sqrt{G \cdot m_e}
	\end{equation}
	
	Aufgelöst nach $G$:
	\begin{equation}
		G_{\text{nat}} = \frac{\xi^2}{4m_e}
	\end{equation}
	
	Dimension: $[G] = [E^{-2}]$ in natürlichen Einheiten; die rechte Seite $\xi^2/(4m_e)$ hat dagegen $[E^{-1}]$ -- der fehlende Faktor $1/E_{\text{char}}$ und die SI-Umrechnung stecken in $C_{\text{conv}}$ (Dok.~012). \textit{(Korrigiert am 30.~Sept.~2026: vorher ohne diesen Hinweis, so dass $\xi^2/(4m_e)$ als $[E^{-2}]$ erschien.)}
	
	% ============================================================================
	\section{Leptonmassen}
	\label{sec:lepton_masses}
	
	Das T0-Modell sagt Leptonmassen voraus (Herleitung in Dokument 003):
	
	\begin{table}[h]
		\centering
		\begin{tabular}{lccc}
			\hline
			\textbf{Lepton} & \textbf{T0 [MeV]} & \textbf{Exp [MeV]} & \textbf{Δ [\%]} \\
			\hline
			Elektron & 0,505 & 0,511 & -1,18 \\
			Myon & 104,96 & 105,66 & -0,66 \\
			Tau & 1783,4 & 1776,9 & +0,37 \\
			\hline
		\end{tabular}
	\end{table}
	\textit{(Korrigiert am 30.~Sept.~2026: vorher $0{,}507/103{,}5/1815$~MeV mit $0{,}87/2{,}09/2{,}16\,\%$ -- die Korpusleiter $m_i = r_i\,\xi^{p_i}\,v$ (Dok.~005, 006) ergibt $0{,}505/104{,}96/1783{,}4$~MeV, also $-1{,}18/-0{,}66/+0{,}37\,\%$; auch $(0{,}511-0{,}507)/0{,}511$ wären $0{,}78\,\%$, nicht $0{,}87\,\%$.)}
	
	% ============================================================================
	\section{Anomale magnetische Momente}
	\label{sec:g2}

	Anomale magnetische Momente werden in eigenen Dokumenten behandelt: Dok.~018 (mit den Revisionen) und Dok.~158.

	\textit{(Korrigiert am 30.~Sept.~2026: vorher stand hier die Kurzformel $a_\ell = \frac{\xi}{2\pi}(E_\ell/E_e)^2$ mit $a_\mu = 1{,}166 \times 10^{-3}$, $a_e = 2{,}122 \times 10^{-5}$ und $a_\tau = 1{,}28 \times 10^{-3}$. Nachgerechnet ergibt die Formel $a_\mu = 2{,}122 \times 10^{-5} \times 42\,753 = 0{,}907$ und $a_\tau = 256{,}6$; $a_e$ liegt um den Faktor 55 unter dem Messwert $1{,}15965 \times 10^{-3}$. Die Formel ist entfernt, ebenso ihre Wiederholungen in Dimensionsanalyse, Formel-Referenz und Berechnungsbeispielen.)}

	% ============================================================================
	\section{Drei Feldgeometrien}
	\label{sec:field_geometries}
	
	\subsection{Typ 1: Lokalisiert sphärisch}
	
	\begin{equation}
		E(r) = E_0 \left(1 - \frac{\beta}{r}\right), \quad \beta = r_0
	\end{equation}
	
	Anwendung: Einzelteilchen (Elektron, Myon, Tau)
	
	\subsection{Typ 2: Lokalisiert nicht-sphärisch}
	
	\begin{equation}
		E(\vec{r}) = E_0 \left(1 - \frac{\beta_{ij} r_i r_j}{r^3}\right)
	\end{equation}
	
	Anwendung: Verbundene Systeme
	
	\subsection{Typ 3: Ausgedehnt homogen}
	
	Effektiver Parameter:
	\begin{equation}
		\xi_{\text{eff}} = \frac{\xi}{2} = \frac{2}{3} \times 10^{-4}
	\end{equation}
	
	Anwendung: Kosmologie (siehe Dokument 026)
	
	% ============================================================================
	\section{Mathematische Identitäten}
	\label{sec:identities}
	
	\subsection{Energiefeld-Normierung}
	
	\begin{equation}
		E_{\text{field}}(\vec{r}, t) = E_0 \cdot f(\vec{r}, t) \cdot e^{i\phi(\vec{r}, t)}
	\end{equation}
	
	mit:
	\begin{itemize}
		\item $E_0$: Charakteristische Energie
		\item $f(\vec{r}, t)$: Normiertes Profil
		\item $\phi(\vec{r}, t)$: Phase
	\end{itemize}
	
	\subsection{Dualitäts-Konsistenz}
	
	Zeit-Masse (Dokument 003): $T \cdot m = 1$
	
	Zeit-Energie (dieses Dokument): $T \cdot E = 1$
	
	In natürlichen Einheiten ($c = 1$):
	\begin{equation}
		E = mc^2 = m \quad \Rightarrow \quad T \cdot m = T \cdot E
	\end{equation}
	
	% ============================================================================
	\section{Dimensionsanalyse-Verifikationen}
	\label{sec:dimensional_analysis}
	
	\subsection{Feldgleichung}
	
	\begin{align}
		[\nabla^2 E] &= [L^{-2}][E] = [E^2][E] = [E^3] \\
		[4\pi G \rho E] &= [E^{-2}][E^4][E] = [E^3] \quad \checkmark
	\end{align}
	
	\subsection{Charakteristische Länge}
	
	\begin{equation}
		[r_0] = [2GE] = [E^{-2}][E] = [E^{-1}] = [L] \quad \checkmark
	\end{equation}
	
	\subsection{Lagrange-Dichte}
	
	\begin{equation}
		[\mathcal{L}] = [\xi][(\partial E)^2] = [1][E^4] = [E^4] \quad \text{(korrekt für Lagrange-Dichte)}
	\end{equation}
	\textit{(Korrigiert am 30.~Sept.~2026: vorher $[1][E^2] = [E^2]$; $[\partial_\mu E] = E^2$, also $[(\partial E)^2] = E^4$, wie im Symbolverzeichnis.)}
	
	% ============================================================================
	\section{Formeln-Referenz}
	\label{sec:formula_reference}
	
	\subsection{Fundamentale Gleichungen}
	
	\begin{align}
		\text{Dualität:} \quad & T_{\text{field}} \cdot E_{\text{field}} = 1 \\
		\text{Wellengleichung:} \quad & \square E_{\text{field}} = 0 \\
		\text{Mit Quellen:} \quad & \nabla^2 E = 4\pi G \rho E \\
		\text{Lagrange-Dichte:} \quad & \mathcal{L} = \xi (\partial E)^2
	\end{align}
	
	\subsection{Abgeleitete Konstanten}
	
	\begin{align}
		\text{Charakteristische Energie:} \quad & E_0 = \sqrt{m_e \cdot m_\mu} = 7{,}348 \text{ MeV (nackt; korrigiert } 7{,}398\text{)} \\
		\text{Feinstrukturkonstante:} \quad & \alpha = \xi (E_0/1\,\text{MeV})^2 \approx 1/137 \\
		\text{Gravitationskonstante:} \quad & G = \frac{\xi^2}{4m_e} \times \text{Faktoren}
	\end{align}
	
	\subsection{Charakteristische Skalen}
	
	\begin{align}
		\text{T0-Länge:} \quad & r_0 = 2GE \\
		\text{T0-Zeit:} \quad & t_0 = 2GE \\
		\text{Planck-Länge:} \quad & \ell_P = \sqrt{G} = 1 \\
		\text{Skalenverhältnis:} \quad & \xi_{\text{ratio}} = \frac{1}{2\sqrt{G} E}
	\end{align}
	
	\subsection{Vorhersageformeln}
	
	\begin{align}
		\text{Parameter:} \quad & \xi = \frac{4}{3} \times 10^{-4} \\
		\text{Effektiver Parameter:} \quad & \xi_{\text{eff}} = \frac{\xi}{2}
	\end{align}
	
	% ============================================================================
	\section{Numerische Werte}
	\label{sec:numerical_values}
	
	\subsection{Fundamentale Konstanten (in natürlichen Einheiten)}
	
	\begin{align}
		\hbar &= 1 \\
		c &= 1 \\
		\alpha &= \frac{1}{137{,}036} \approx 7{,}297 \times 10^{-3} \\
		G &= 1 \text{ (numerisch, Dimension } [E^{-2}]\text{)}
	\end{align}
	
	\subsection{T0-Parameter}
	
	\begin{align}
		\xi &= \frac{4}{3} \times 10^{-4} = 1{,}3333 \times 10^{-4} \\
		\xi^2 &= 1{,}7778 \times 10^{-8} \\
		\frac{\xi}{2\pi} &= 2{,}1221 \times 10^{-5} \\
		\xi_{\text{eff}} &= 6{,}6667 \times 10^{-5} \\
		E_0 &= 7{,}348 \text{ MeV (aus exp. Massen)} \\
		E_0^{\text{T0}} &= 7{,}398 \text{ MeV (fraktal korrigiert, } E_0/\sqrt{K_{\text{frak}}}\text{)}
	\end{align}
	
	\subsection{Leptonenergieen}
	
	\begin{align}
		E_e &= 0{,}511 \text{ MeV} \\
		E_\mu &= 105{,}658 \text{ MeV} \\
		E_\tau &= 1776{,}86 \text{ MeV}
	\end{align}
	
	\subsection{Energieverhältnisse}
	
	\begin{align}
		\frac{E_\mu}{E_e} &= 206{,}768 \\
		\frac{E_\tau}{E_e} &= 3477{,}2 \\
		\frac{E_\tau}{E_\mu} &= 16{,}817
	\end{align}
	
	% ============================================================================
	\section{Berechnungsbeispiele}
	\label{sec:calculations}
	
	\subsection{Feinstrukturkonstante}
	
	Gegeben:
	\begin{itemize}
		\item $\xi = 1{,}3333 \times 10^{-4}$
		\item $E_0 = 7{,}398$ MeV
	\end{itemize}
	
	Berechnung:
	\begin{align}
		\left(\frac{E_0}{1\,\text{MeV}}\right)^2 &= (7{,}398)^2 = 54{,}73 \\
		\alpha &= 1{,}3333 \times 10^{-4} \times 54{,}73 \\
		&= 7{,}297 \times 10^{-3} \\
		&= \frac{1}{137{,}04}
	\end{align}
	
	Experimentell: $\alpha_{\text{exp}} = \frac{1}{137{,}036}$
	
	Abweichung: 0,0005\,\% \textit{(korrigiert am 30.~Sept.~2026: vorher 0,03\,\%)}
	
	\subsection{Charakteristische Länge (Elektron)}
	
	Gegeben:
	\begin{itemize}
		\item $E_e = 0{,}511$ MeV $= 0{,}511 \times 1{,}6 \times 10^{-13}$ J $= 8{,}2 \times 10^{-14}$ J
		\item $G = 6{,}674 \times 10^{-11}$ m$^3$ kg$^{-1}$ s$^{-2}$
		\item $c = 3 \times 10^8$ m/s
	\end{itemize}
	
	Umrechnung in SI-Einheiten:
	\begin{equation}
		r_0 = \frac{2GE}{c^4} \approx 1{,}35 \times 10^{-57} \text{ m}
	\end{equation}
	
	Planck-Vergleich:
	\begin{equation}
		\frac{r_0}{\ell_P} = \frac{1{,}35 \times 10^{-57}}{1{,}6 \times 10^{-35}} \approx 8{,}4 \times 10^{-23}
	\end{equation}
	
	\textit{(Korrigiert am 29.~Sept.~2026: vorher $r_0 = 2GE \approx 10^{-28}$~m und $r_0/\ell_P \approx 10^7$ -- mit den angegebenen SI-Werten lautet die Formel $r_0 = 2GE/c^4$ und ergibt $1{,}35 \times 10^{-57}$~m; $r_0$ liegt damit weit unterhalb der Planck-Länge.)}
	
	% ============================================================================
	% APPENDIX
	% ============================================================================
	
	\appendix
	
	\section{Symbolverzeichnis}
	
	\begin{longtable}{|c|l|c|}
		\hline
		\textbf{Symbol} & \textbf{Bedeutung} & \textbf{Dimension} \\
		\hline
		$\xi$ & Fundamentaler Parameter & $1$ \\
		$E_0$ & Charakteristische Energie & $E$ \\
		$E_{\text{field}}$ & Universelles Energiefeld & $E$ \\
		$T_{\text{field}}$ & Intrinsisches Zeitfeld & $E^{-1}$ \\
		$r_0$ & T0-charakteristische Länge & $L = E^{-1}$ \\
		$t_0$ & T0-charakteristische Zeit & $T = E^{-1}$ \\
		$\ell_P$ & Planck-Länge & $L = E^{-1}$ \\
		$G$ & Gravitationskonstante & $E^{-2}$ \\
		$\alpha$ & Feinstrukturkonstante & $1$ \\
		$E_e, E_\mu, E_\tau$ & Leptonenergieen & $E$ \\
		$m_e, m_\mu, m_\tau$ & Leptonmassen ($= E$ in nat. Einh.) & $E$ \\
		$\mathcal{L}$ & Lagrange-Dichte & $E^4$ \\
		$\square$ & d'Alembert-Operator & $E^2$ \\
		$\xi_{\text{eff}}$ & Effektiver Parameter ($\xi/2$) & $1$ \\
		\hline
	\end{longtable}
	
	\section{Einheiten-Umrechnungen}
	
	\subsection{Natürliche → SI}
	
	\begin{align}
		1 \text{ (Energie)} &= 1 \text{ GeV} = 1{,}6 \times 10^{-10} \text{ J} \\
		1 \text{ (Länge)} &= \frac{\hbar c}{1 \text{ GeV}} = 0{,}197 \text{ fm} \\
		1 \text{ (Zeit)} &= \frac{\hbar}{1 \text{ GeV}} = 6{,}58 \times 10^{-25} \text{ s}
	\end{align}
	
	\subsection{Standard natürliche Einheiten}
	
	In Standard-Konvention ($\hbar = c = 1$):
	\begin{itemize}
		\item $\alpha = \frac{e^2}{4\pi\epsilon_0} \approx \frac{1}{137}$ (dimensionslos)
		\item Alle Größen in Potenzen von Energie
		\item Physikalische Vorhersagen identisch zu anderen Konventionen
	\end{itemize}
	
	\section{Beziehung zu anderen Dokumenten}
	
	\begin{itemize}
		\item \textbf{Dokument 003}: Zeit-Masse-Dualität, Grundlagen, Ursprung von $\xi$
		\item \textbf{Dokument 018}: Geometrische g-2-Formulierung (fraktale Geometrie)
		\item \textbf{Dokument 019}: Lagrangian-Formulierung (Quantenfeldtheorie)
		\item \textbf{Dokument 026}: Kosmologie ($\xi_{\text{eff}} = \xi/2$)
	\end{itemize}
	
	Alle Formulierungen basieren auf $\xi = \frac{4}{3} \times 10^{-4}$.
	